% ======================================================================
\section{Rapidity logarithms: BFKL evolution of the projectile and JIMWLK evolution of the target}\label{sec:rapidity}
% ======================================================================

In this section we collect all terms of the $O(g^4)$ cross section that are enhanced by a logarithm of the longitudinal cut-offs and show that they are the leading-logarithmic evolution of the LO spectrum: the BFKL evolution of the projectile charge correlator over the rapidity interval between the projectile and the produced gluon, and the JIMWLK evolution of the target operator over the interval between the produced gluon and the target.

\subsection{Origin of the logarithms}\label{sec:loglong}

Rapidity logarithms arise from longitudinal integrals of the form $\int dp^+/p^+$ over the momentum of the second gluon, real (integrated over in Groups~II and~III) or virtual (inside the coefficients $\mf N_2$, $\mf A^{(3)}$, $\mf B_3$ and in the interference terms of Group~I). In every row the integrand behaves as $1/p^+$ in two regions: for $k^+\ll p^+<\vee$, where the second gluon is faster than the observed one and belongs to the projectile, and for $\Lambda<p^+\ll k^+$, where it is slower and is resolved by the observed gluon as part of the target. The first region produces $\log(\vee/k^+)$, the second $\log(k^+/\Lambda)$. Terms that do not depend on $k^+$ at all, in particular the normalisation $\mf N_2$ and the products of four single-emission amplitudes, carry the full logarithm
\begin{align}\label{eq:logsplit}
\log\frac{\vee}{\Lambda}=\log\frac{\vee}{k^+}+\log\frac{k^+}{\Lambda}\,,
\end{align}
which we split at the longitudinal momentum $k^+$ of the produced gluon. With this split every logarithmic term is assigned either to the projectile or to the target side. Terms proportional to $\log(1-k^+/\vee)$, which arise from the upper limit $\vee-k^+$ of the integrals with a soft parent gluon, vanish for $\vee\to\infty$ and are dropped.

By the order of the charges, the logarithmic terms are quartic, cubic or quadratic in $\rho$ (Table~\ref{tab:rows}). We discuss them in this order, because the reordering of the non-commuting charges feeds terms with more charges into terms with fewer.

\subsection{Terms quartic in \texorpdfstring{$\rho$}{rho}}\label{sec:fourrho}

Terms with four charges arise in Rows I.I, I.II (through the three-$\rho$ part $\mf A\mf N_2$ of $\mf A^{(3)}$), I.IV, II.I, II.II, II.III, III.I, III.II and III.IV, from the products of four single-emission coefficients $\mf A$, of $\mf A$ with $\mf N_2$, and of $\mf A$ with the $\{\rho,\rho\}$ part of $\mf B_2$. In all of them the longitudinal integral runs over the full interval $\Lambda<p^+<\vee$ and gives $\log(\vee/\Lambda)$, split according to Eq.~\eqref{eq:logsplit}.

The totally symmetrised four-$\rho$ part of their sum vanishes identically, i.e.\ the four-$\rho$ terms cancel for commuting charges. This cancellation has a simple origin. For commuting charges the part of $\Omega$ that is linear in $\rho$ exponentiates: to the order considered, $\Omega$ acting on the vacuum is the coherent state
\begin{align}\label{eq:coherent}
\Omega_{\rm coh}\ket0=\exp\Big\{\int\frac{d^3p}{(2\pi)^3}\Big[\mf A^b_j(p)\,a^{\dagger b}_j(p)-\mf A^{b\dagger}_j(p)\,a^b_j(p)\Big]\Big\}\ket0\,,
\end{align}
whose expansion to second order produces $\mf N_2$ and the $\{\rho,\rho\}$ part of $\mf B_2$, $\mf B_2\supset\frac12\mf A\mf A$. Since $\bar\Omega_{\rm coh}=S^\dagger\Omega_{\rm coh}S$ is again a displacement operator, $\bar\Omega^\dagger_{\rm coh}\Omega_{\rm coh}$ is, up to a phase, a displacement operator, and the expectation value of the number operator in the corresponding coherent state is the squared modulus of the displacement. The displacement is the LO amplitude~\eqref{eq:ALO}. Hence for commuting charges the spectrum is exactly the LO result, with no correction of order $g^4\rho^4$. Two features of the coefficients are necessary for this. First, the $\{\rho,\rho\}$ part of $\mf B_2$ must have support on the full square $\Lambda<k^+,p^+<\vee$, as follows from the kinematics of two independent emissions; with the support $k^++p^+<\vee$ the four-$\rho$ terms of the different rows would carry different logarithms and would not cancel. Second, the three-$\rho$ part of $\mf A^{(3)}$ must be ordered as $\mf A\mf N_2$ (Section~\ref{sec:coefficients}), so that the normalisation terms combine into Eq.~\eqref{eq:Nterms}. We have verified the cancellation by an explicit term-by-term evaluation, with the charges represented by $SU(2)$ and $SU(3)$ matrices and the order of all operators kept.

For non-commuting charges the four-$\rho$ terms do not vanish as operators: rewriting each product in terms of totally symmetrised products leaves remainders with commutators. For two and three operators
\begin{align}\label{eq:symmetrise}
\rho_1\rho_2&=\tfrac12\{\rho_1,\rho_2\}+\tfrac12[\rho_1,\rho_2]\,,\notag\\
\rho_1\rho_2\rho_3&=S(\rho_1\rho_2\rho_3)+\tfrac12\Big(S\big([\rho_1,\rho_2]\rho_3\big)+S\big([\rho_1,\rho_3]\rho_2\big)+S\big([\rho_2,\rho_3]\rho_1\big)\Big)\notag\\&\quad+\tfrac16\Big(\big[[\rho_1,\rho_2],\rho_3\big]-\big[[\rho_2,\rho_3],\rho_1\big]\Big)\,,
\end{align}
where $S$ denotes the totally symmetrised product, $S(XY)=\frac12\{X,Y\}$, and similarly for four operators. By Eq.~\eqref{eq:rho} each commutator reduces the number of charges by one. The four-$\rho$ terms therefore leave a three-$\rho$ remainder, multiplied by $\log(\vee/k^+)$ and by $\log(k^+/\Lambda)$, and a two-$\rho$ remainder, which is part of the BFKL and JIMWLK terms below.

\subsection{Terms cubic in \texorpdfstring{$\rho$}{rho}}\label{sec:threerho}

Logarithmic terms with three charges arise directly in Rows I.II--I.V, II.I, III.I--III.III, III.V and III.VI, from the parts of $\mf B_2$, $\mf B_3$ and $\mf A^{(3)}$ that are linear or quadratic in $\rho$ multiplied by single emissions, and in addition from the reordering of the four-$\rho$ terms. They are of the form $if^{abc}\times$(real kernel)$\times$(Wilson lines)$\times\rho\rho\rho$, they change sign together with the sign of the exponent $e^{-i\v k\cdot(\v w-\v w')}$, and they are of the type of the ``first saturation corrections'' in the projectile.

For each of the two logarithms, the totally symmetrised three-$\rho$ part of the sum of all these terms vanishes identically, for an arbitrary function of $\v w$ and $\v w'$ in place of the phase. We have organised the sum by transverse kernel and by the number of Wilson lines (two, three and four) and verified the cancellation in each class separately, analytically, at the operator level with explicit $SU(2)$ matrices for the charges on a transverse lattice, so that contact terms are included, and for the totally symmetrised colour tensors with random $SU(3)$ and $SU(4)$ Wilson lines (to $2\times10^{-15}$). Since the colour tensor itself vanishes, the cancellation holds for any symmetrised three-point function of the projectile, in particular for the $d^{abc}$ (odderon) one of Eq.~\eqref{eq:A2d}, whose projection we checked separately: no rapidity logarithm multiplies the three-point function. That no logarithm can multiply $\sigma$ follows from charge conjugation. C acts as $\rho^at^a\to-(\rho^at^a)^T$, i.e.\ $\rho^a\to-\eta_a\rho^a$ with $(t^a)^T=\eta_at^a$, on the soft gluons as $A^a\to-\eta_aA^a$ and on the target as $V\to V^*$, $U^{ab}\to\eta_a\eta_bU^{ab}$; since $\eta_a\eta_b\eta_cf^{abc}=-f^{abc}$ and $\eta_a\eta_b\eta_cd^{abc}=d^{abc}$, it preserves Eq.~\eqref{eq:rho} and all couplings, and the spectrum is invariant under the joint transformation. Under C, $n$ is even and $\sigma$ odd, and the Wilson-line traces with $T^e$ are even and those with $D^h$ odd (Section~\ref{sec:fin_singlet}). By the factorisation of Section~\ref{sec:proofsketch}, the coefficient of $\log(\vee/k^+)$ is the LO spectrum with the projectile correlator replaced by its one-step evolution, which does not involve the target: a three-charge term can enter it only through a $c$-number invariant tensor, $d^{abc}$ or $f^{abc}$, and C-invariance of the evolution admits only $f^{abc}$, whose contraction with Eq.~\eqref{eq:A2d} vanishes; the target enters through the C-even $\mathcal O_{\rm LO}$. The coefficient of $\log(k^+/\Lambda)$ is $-H_{\rm JIMWLK}$ acting on the target operator of the LO spectrum, in which the projectile enters only through $n$. Hence neither logarithm can multiply $\sigma$; a dependence on $\sigma$ requires a C-odd target structure and is non-logarithmic (Section~\ref{sec:fin_rows}). C alone does not forbid terms $\sigma\times$(C-odd target structure), which are exactly the symmetrised three-$\rho$ finite terms; the factorised form of the logarithms is needed in addition. The explicit cancellation confirms this argument and is stronger: individual terms have $d^{abc}$ projections of order one, with C-odd target factors, which C alone would allow, and they cancel in the sum; the cancellation of the whole symmetrised tensor also removes an $f^{abc}g_f$ part, which is C-even and not forbidden by C. As operators the three-$\rho$ terms leave, by Eq.~\eqref{eq:symmetrise}, two-$\rho$ and one-$\rho$ remainders. The two-$\rho$ remainders are part of the BFKL and JIMWLK terms. The one-$\rho$ remainders are contact terms $\propto f^{abc}\rho^c(\v x)\delta^{(2)}(\v x-\v y)$, which do not contribute in the colour-singlet form.

The non-logarithmic three-$\rho$ terms do not cancel. They belong to the finite terms (Section~\ref{sec:finite}).

\subsection{The BFKL term}\label{sec:BFKL}

Collecting all two-$\rho$ terms proportional to $\log(\vee/k^+)$, the genuine ones of the rows of Table~\ref{tab:rows} (I.II--I.IV, II.I, II.IV, III.I--III.III, III.V, III.VI) and the remainders of the reordered three- and four-$\rho$ terms, we find
\begin{align}\label{eq:BFKLrho}
\log\frac{\vee}{k^+}\,\frac{d^{3}N_{\rm BFKL}}{d^{3}k} & = \frac{1}{(2\pi)^{3}}\,\frac{g^{4}}{16\pi^{5}}\,\frac{1}{k^{+}}\,\log\Big[\frac{\vee}{k^{+}}\Big]\notag\\&\times\int_{x,x',z,w,w'}e^{-ik\cdot (w-w')}\Bigg[ \notag\\
& +\mathcal{K}(x-w)\cdot\mathcal{K}(x'-w')\,\frac{1}{(x-z)^{2}}\Big[\notag\\
& -N_{c}\,U^{ab}(w')\,U^{ac}(w)\,\rho^{b}(x')\,\rho^{c}(x)\notag\\
&\quad +N_{c}\,U^{ab}(w')\,U^{ac}(x)\,\rho^{b}(x')\,\rho^{c}(x)\notag\\
& +N_{c}\,U^{ab}(x')\,U^{ac}(w)\,\rho^{b}(x')\,\rho^{c}(x) \notag\\
&\quad -N_{c}\,U^{ab}(x')\,U^{ac}(x)\,\rho^{b}(x')\,\rho^{c}(x)\Big]\notag\\
& +\mathcal{K}(x-w)\cdot\mathcal{K}(x'-w')\,\frac{1}{(x'-z)^{2}}\Big[\notag\\
& -N_{c}\,U^{ab}(w')\,U^{ac}(w)\,\rho^{b}(x')\,\rho^{c}(x)\notag\\
&\quad +N_{c}\,U^{ab}(w')\,U^{ac}(x)\,\rho^{b}(x')\,\rho^{c}(x)\notag\\
& +N_{c}\,U^{ab}(x')\,U^{ac}(w)\,\rho^{b}(x')\,\rho^{c}(x) \notag\\
&\quad -N_{c}\,U^{ab}(x')\,U^{ac}(x)\,\rho^{b}(x')\,\rho^{c}(x)\Big]\notag\\
& +\mathcal{K}(x-w)\cdot\mathcal{K}(x'-w')\,\mathcal{K}(x-z)\cdot\mathcal{K}(x'-z)\Big[\notag\\
& +2\,f^{abc}\,f^{dbe}\,U^{ha}(w')\,U^{hd}(w)\,\rho^{c}(x')\,\rho^{e}(x)\notag\\
&\quad -2\,f^{abc}\,f^{dbe}\,U^{ha}(w')\,U^{hd}(x)\,\rho^{c}(x')\,\rho^{e}(x)\notag\\
& -2\,f^{abc}\,f^{dbe}\,U^{ha}(x')\,U^{hd}(w)\,\rho^{c}(x')\,\rho^{e}(x) \notag\\
&\quad +2\,f^{abc}\,f^{dbe}\,U^{ha}(x')\,U^{hd}(x)\,\rho^{c}(x')\,\rho^{e}(x)\Big]\notag\\
& +\mathcal{K}(w-z)\cdot\mathcal{K}(w'-z)\,\mathcal{K}(x-z)\cdot\mathcal{K}(x'-z)\Big[\notag\\
& +2\,N_{c}\,\rho^{a}(x)\,\rho^{a}(x')\notag\\
&\quad -\,f^{abc}\,f^{dce}\,U^{ha}(w')\,U^{hd}(w)\,\rho^{b}(x)\,\rho^{e}(x')\notag\\
& -\,f^{abc}\,f^{dce}\,U^{ha}(w')\,U^{hd}(w)\,\rho^{b}(x')\,\rho^{e}(x) \notag\\
& +\,f^{abc}\,f^{dce}\,U^{ha}(w')\,U^{hd}(z)\,\rho^{b}(x)\,\rho^{e}(x')\notag\\
&\quad +\,f^{abc}\,f^{dce}\,U^{ha}(w')\,U^{hd}(z)\,\rho^{b}(x')\,\rho^{e}(x)\notag\\
& +\,f^{abc}\,f^{dce}\,U^{ha}(z)\,U^{hd}(w)\,\rho^{b}(x)\,\rho^{e}(x') \notag\\
&\quad +\,f^{abc}\,f^{dce}\,U^{ha}(z)\,U^{hd}(w)\,\rho^{b}(x')\,\rho^{e}(x)\Big]\notag\\
& +\mathcal{K}(x'-w')\cdot\mathcal{K}(w-z)\,\frac{1}{(x-z)^{2}}\Big[\notag\\
& -N_{c}\,U^{ab}(w')\,U^{ac}(w)\,\rho^{b}(x')\,\rho^{c}(x) \notag\\
& +N_{c}\,U^{ab}(w')\,U^{ac}(z)\,\rho^{b}(x')\,\rho^{c}(x)\notag\\
&\quad +N_{c}\,U^{ab}(x')\,U^{ac}(w)\,\rho^{b}(x')\,\rho^{c}(x)\notag\\
& -N_{c}\,U^{ab}(x')\,U^{ac}(z)\,\rho^{b}(x')\,\rho^{c}(x)\Big]\notag\\
& +\mathcal{K}(x'-w')\cdot\mathcal{K}(w-z)\,\mathcal{K}(x-z)\cdot\mathcal{K}(x'-z)\Big[\notag\\
& +2\,f^{abc}\,f^{cde}\,U^{hd}(w')\,U^{ha}(w)\,\rho^{b}(x)\,\rho^{e}(x')\notag\\
&\quad -2\,f^{abc}\,f^{cde}\,U^{hd}(w')\,U^{ha}(z)\,\rho^{b}(x)\,\rho^{e}(x')\notag\\
& -2\,f^{abc}\,f^{cde}\,U^{hd}(x')\,U^{ha}(w)\,\rho^{b}(x)\,\rho^{e}(x') \notag\\
&\quad +2\,f^{abc}\,f^{cde}\,U^{hd}(x')\,U^{ha}(z)\,\rho^{b}(x)\,\rho^{e}(x')\Big]\notag\\
& +\mathcal{K}(x-w)\cdot\mathcal{K}(w'-z)\,\frac{1}{(x'-z)^{2}}\Big[\notag\\
& -N_{c}\,U^{ab}(w')\,U^{ac}(w)\,\rho^{c}(x)\,\rho^{b}(x')\notag\\
&\quad +N_{c}\,U^{ab}(w')\,U^{ac}(x)\,\rho^{c}(x)\,\rho^{b}(x')\notag\\
& +N_{c}\,U^{ab}(z)\,U^{ac}(w)\,\rho^{c}(x)\,\rho^{b}(x') \notag\\
&\quad -N_{c}\,U^{ab}(z)\,U^{ac}(x)\,\rho^{c}(x)\,\rho^{b}(x')\Big]\notag\\
& +\mathcal{K}(x-w)\cdot\mathcal{K}(w'-z)\,\mathcal{K}(x-z)\cdot\mathcal{K}(x'-z)\Big[\notag\\
& +2\,f^{abc}\,f^{cde}\,U^{ha}(w')\,U^{hd}(w)\,\rho^{b}(x')\,\rho^{e}(x)\notag\\&\quad -2\,f^{abc}\,f^{cde}\,U^{ha}(w')\,U^{hd}(x)\,\rho^{b}(x')\,\rho^{e}(x)\notag\\
& -2\,f^{abc}\,f^{cde}\,U^{ha}(z)\,U^{hd}(w)\,\rho^{b}(x')\,\rho^{e}(x) \notag\\
&\quad +2\,f^{abc}\,f^{cde}\,U^{ha}(z)\,U^{hd}(x)\,\rho^{b}(x')\,\rho^{e}(x)\Big]\Bigg]
\end{align}
where the average over the projectile is implied. With non-commuting charges the projectile evolution produces the symmetrised product $\frac12\{\rho^b(\v x'),\rho^c(\v x)\}$ in the last four blocks; the fixed order written there differs from it by contact terms $\propto f^{abc}\rho^c(\v x)\delta^{(2)}(\v x-\v x')$, which do not contribute in the colour-singlet form. The prefactor is $\frac{1}{(2\pi)^3}\frac{g^4}{16\pi^5}=\frac{\alpha_s^2}{8\pi^6}$.

For a colour-singlet projectile, Eq.~\eqref{eq:singlet}, Eq.~\eqref{eq:BFKLrho} becomes
\begin{align}\label{eq:BFKLsinglet}
\log\frac{\vee}{k^+}\,\frac{d^{3}N_{\rm BFKL}}{d^{3}k} & = \frac{1}{(2\pi)^{3}}\,\frac{g^{4}}{16\pi^{5}}\,\frac{1}{k^{+}}\,\frac{1}{N^{2}_{c}-1}\,\log\Big[\frac{\vee}{k^{+}}\Big]\,\int_{x,y,z,w,w'}e^{-ik\cdot (w-w')}\notag\\&\times\langle\rho^{f}(x)\rho^{f}(y) \rangle\,\Bigg[ \notag\\ & \mathcal K(x-w)\cdot\mathcal K(y-w')\,K(x,y,z)\Big[2 N_{c}\mathcal O_{\rm LO}(x,y;w,w')\Big]\notag\\
&+\mathcal K(x-w)\cdot\mathcal K(y-w')\,K(x,x,z)\Big[-N_{c}\mathcal O_{\rm LO}(x,y;w,w')\Big]\notag\\
&+\mathcal K(x-w)\cdot\mathcal K(y-w')\,K(y,y,z)\Big[-N_{c}\mathcal O_{\rm LO}(x,y;w,w')\Big]\notag\\
&+K(w,w',z)\,K(x,y,z)\Big[2 N_{c}\mathcal O_{\rm LO}(z,z;w,w')\Big]\notag\\
&+\mathcal K(x-w')\cdot\mathcal K(w-z)\,K(x,y,z)\Big[2 N_{c}\mathcal O_{\rm LO}(z,x;w,w')\Big]\notag\\
&+\mathcal K(x-w')\cdot\mathcal K(w-z)\,K(y,y,z)\Big[-N_{c}\mathcal O_{\rm LO}(z,x;w,w')\Big]\notag\\
&+\mathcal K(x-w)\cdot\mathcal K(w'-z)\,K(x,y,z)\Big[2 N_{c}\mathcal O_{\rm LO}(x,z;w,w')\Big]\notag\\
&+\mathcal K(x-w)\cdot\mathcal K(w'-z)\,K(y,y,z)\Big[-N_{c}\mathcal O_{\rm LO}(x,z;w,w')\Big]\Bigg]
\end{align}
with $K(\v x,\v y,\v z)$ of Eq.~\eqref{eq:kernels} and $\mathcal O_{\rm LO}$ of Eq.~\eqref{eq:OLO}. As for LO, the sign of the exponent is immaterial in this form.

Eq.~\eqref{eq:BFKLsinglet} is the LO spectrum~\eqref{eq:LOsinglet} with the charge correlator replaced by its one-step BFKL evolution over the rapidity interval $\log(\vee/k^+)$:
\begin{align}\label{eq:BFKLasLO}
\log\frac{\vee}{k^+}\,\frac{d^3N_{\rm BFKL}}{d^3k}=\log\frac{\vee}{k^+}\;\frac{d^3N_{\rm LO}}{d^3k}\bigg|_{n(\v x,\v y)\to\frac{d}{dY}n(\v x,\v y)}\,,
\end{align}
where the evolution of the projectile charge correlator in the dilute regime is~\cite{Kovner:2006wr,Kovner:2006ge}
\begin{align}\label{eq:BFKLrhorho}
\frac{d}{dY}\big\langle\rho^f(\v x)\rho^f(\v y)\big\rangle=-\frac{\alpha_sN_c}{2\pi^2}\int_z\Big[&K_D(\v x,\v y,\v z)\,\big\langle\rho^f(\v x)\rho^f(\v y)\big\rangle
\notag\\&-2\,\delta^{(2)}(\v x-\v y)\int_{z'}K(\v z,\v z',\v x)\,\big\langle\rho^f(\v z)\rho^f(\v z')\big\rangle
\notag\\&
+\big(2K(\v y,\v z,\v x)-K(\v z,\v z,\v x)\big)\big\langle\rho^f(\v y)\rho^f(\v z)\big\rangle
\notag\\&+\big(2K(\v x,\v z,\v y)-K(\v z,\v z,\v y)\big)\big\langle\rho^f(\v x)\rho^f(\v z)\big\rangle\Big]\,,
\end{align}
with the dipole kernel $K_D(\v x,\v y,\v z)=K(\v x,\v x,\v z)-2K(\v x,\v y,\v z)+K(\v y,\v y,\v z)$. The first term of Eq.~\eqref{eq:BFKLrhorho} gives the first three lines of Eq.~\eqref{eq:BFKLsinglet} (virtual terms, the produced gluon is emitted by the original charges); the second term, in which the charge correlator is concentrated at the position $\v z$ of the new gluon, gives the fourth line (the produced gluon is emitted by the new gluon); and the last two terms give the remaining four lines (the produced gluon is emitted by the new gluon in the amplitude and by an original charge in the conjugate amplitude, or vice versa). In terms of the colour-singlet correlator of the dipole picture, Eq.~\eqref{eq:BFKLrhorho} is the linearised BK (dipole BFKL) equation; the check for a single colour dipole is described in Section~\ref{sec:checks}.

\subsection{The JIMWLK term}\label{sec:JIMWLK}

The two-$\rho$ terms proportional to $\log(k^+/\Lambda)$, the genuine ones of Rows I.II, I.V, II.I, II.IV and III.III and the remainders of the reordered three- and four-$\rho$ terms, add up to
\begin{align}\label{eq:JIMWLKrho}
\log\frac{k^+}{\Lambda}\,\frac{d^{3}N_{\rm JIMWLK}}{d^{3}k} & = \frac{1}{(2\pi)^{3}}\,\frac{g^{4}}{16\pi^{5}}\,\frac{1}{k^{+}}\,\log{\left(\frac{k^{+}}{\Lambda}\right)}\,\int_{x,x',z,w,w'}e^{-ik\cdot(w-w')}\,\Bigg[
\notag \\ & \mathcal{K}(x-w)\cdot\mathcal{K}(x'-w')\,\frac{1}{(x-z)^{2}}\Big[\notag\\
&+2\,N_{c}\,U^{ab}(w')\,U^{ac}(x)\,\rho^{b}(x')\,\rho^{c}(x) -2\,N_{c}\,U^{ab}(x')\,U^{ac}(x)\,\rho^{b}(x')\,\rho^{c}(x)\notag\\
& +2\,f^{abc}\,f^{deh}\,U^{da}(z)\,U^{el}(w')\,U^{hb}(x)\,\rho^{l}(x')\,\rho^{c}(x) \notag\\
&\quad -2\,f^{abc}\,f^{deh}\,U^{da}(z)\,U^{el}(x')\,U^{hb}(x)\,\rho^{l}(x')\,\rho^{c}(x)\Big]\notag\\
& +\mathcal{K}(x-w)\cdot\mathcal{K}(x'-w')\,\frac{1}{(x'-z)^{2}}\Big[\notag\\
& +2\,N_{c}\,U^{ab}(x')\,U^{ac}(w)\,\rho^{b}(x')\,\rho^{c}(x)\notag\\
&\quad -2\,N_{c}\,U^{ab}(x')\,U^{ac}(x)\,\rho^{b}(x')\,\rho^{c}(x)\notag\\
& +2\,f^{abc}\,f^{deh}\,U^{da}(z)\,U^{el}(w)\,U^{hb}(x')\,\rho^{c}(x')\,\rho^{l}(x) \notag\\
&\quad -2\,f^{abc}\,f^{deh}\,U^{da}(z)\,U^{el}(x)\,U^{hb}(x')\,\rho^{c}(x')\,\rho^{l}(x)\Big]\notag\\
& +\mathcal{K}(x-w)\cdot\mathcal{K}(x'-w')\,\mathcal{K}(x-z)\cdot\mathcal{K}(x'-z)\Big[\notag\\
& +2\,N_{c}\,U^{ab}(x')\,U^{ac}(x)\,\rho^{b}(x')\,\rho^{c}(x)\notag\\
&\quad +2\,f^{abc}\,f^{ade}\,U^{hb}(x')\,U^{hd}(x)\,\rho^{c}(x')\,\rho^{e}(x)\notag\\
& -2\,f^{abc}\,f^{deh}\,U^{da}(z)\,U^{eb}(x)\,U^{hl}(x')\,\rho^{l}(x')\,\rho^{c}(x) \notag\\
&\quad -2\,f^{abc}\,f^{deh}\,U^{da}(z)\,U^{eb}(x')\,U^{hl}(x)\,\rho^{c}(x')\,\rho^{l}(x)\Big]\notag\\
& +\mathcal{K}(x-w)\cdot\mathcal{K}(x'-w')\,\frac{1}{(w-z)^{2}}\Big[\notag\\
& -2\,N_{c}\,U^{ab}(w')\,U^{ac}(w)\,\rho^{b}(x')\,\rho^{c}(x)\notag\\
&\quad +2\,N_{c}\,U^{ab}(x')\,U^{ac}(w)\,\rho^{b}(x')\,\rho^{c}(x)\notag\\
& -2\,f^{abc}\,f^{deh}\,U^{da}(z)\,U^{el}(w')\,U^{hb}(w)\,\rho^{l}(x')\,\rho^{c}(x) \notag\\
&\quad +2\,f^{abc}\,f^{deh}\,U^{da}(z)\,U^{el}(x')\,U^{hb}(w)\,\rho^{l}(x')\,\rho^{c}(x)\Big]\notag\\
& +\mathcal{K}(x-w)\cdot\mathcal{K}(x'-w')\,\frac{1}{(w'-z)^{2}}\Big[\notag\\
& -2\,N_{c}\,U^{ab}(w')\,U^{ac}(w)\,\rho^{b}(x')\,\rho^{c}(x)\notag\\
&\quad +2\,N_{c}\,U^{ab}(w')\,U^{ac}(x)\,\rho^{b}(x')\,\rho^{c}(x)\notag\\
& -2\,f^{abc}\,f^{deh}\,U^{da}(z)\,U^{el}(w)\,U^{hb}(w')\,\rho^{c}(x')\,\rho^{l}(x) \notag\\
&\quad +2\,f^{abc}\,f^{deh}\,U^{da}(z)\,U^{el}(x)\,U^{hb}(w')\,\rho^{c}(x')\,\rho^{l}(x)\Big]\notag\\
& +\mathcal{K}(x-w)\cdot\mathcal{K}(x'-w')\,\mathcal{K}(w-z)\cdot\mathcal{K}(w'-z)\Big[\notag\\
& +2\,N_{c}\,U^{ab}(w')\,U^{ac}(w)\,\rho^{b}(x')\,\rho^{c}(x)\notag\\
&\quad +2\,f^{abc}\,f^{ade}\,U^{hb}(w')\,U^{hd}(w)\,\rho^{c}(x')\,\rho^{e}(x)\notag\\
& -2\,f^{abc}\,f^{deh}\,U^{da}(z)\,U^{eb}(w)\,U^{hl}(w')\,\rho^{l}(x')\,\rho^{c}(x) \notag\\
&\quad -2\,f^{abc}\,f^{deh}\,U^{da}(z)\,U^{eb}(w')\,U^{hl}(w)\,\rho^{c}(x')\,\rho^{l}(x)\Big]\notag\\
& +\mathcal{K}(x-w)\cdot\mathcal{K}(x'-w')\,\mathcal{K}(x-z)\cdot\mathcal{K}(w'-z)\Big[\notag\\
& -2\,N_{c}\,U^{ab}(w')\,U^{ac}(x)\,\rho^{b}(x')\,\rho^{c}(x)\notag\\
& -2\,f^{abc}\,f^{ade}\,U^{hb}(w')\,U^{hd}(x)\,\rho^{c}(x')\,\rho^{e}(x) \notag\\
&\quad +2\,f^{abc}\,f^{deh}\,U^{da}(z)\,U^{eb}(x)\,U^{hl}(w')\,\rho^{l}(x')\,\rho^{c}(x)\notag\\
& +2\,f^{abc}\,f^{deh}\,U^{da}(z)\,U^{eb}(w')\,U^{hl}(x)\,\rho^{c}(x')\,\rho^{l}(x)\Big]\notag\\
& +\mathcal{K}(x-w)\cdot\mathcal{K}(x'-w')\,\mathcal{K}(x'-z)\cdot\mathcal{K}(w-z)\Big[\notag\\
& -2\,N_{c}\,U^{ab}(x')\,U^{ac}(w)\,\rho^{b}(x')\,\rho^{c}(x)\notag\\&\quad -2\,f^{abc}\,f^{ade}\,U^{hb}(x')\,U^{hd}(w)\,\rho^{c}(x')\,\rho^{e}(x)\notag\\
& +2\,f^{abc}\,f^{deh}\,U^{da}(z)\,U^{eb}(w)\,U^{hl}(x')\,\rho^{l}(x')\,\rho^{c}(x) \notag\\
&\quad +2\,f^{abc}\,f^{deh}\,U^{da}(z)\,U^{eb}(x')\,U^{hl}(w)\,\rho^{c}(x')\,\rho^{l}(x)\Big]\Bigg]
\end{align}
For a colour-singlet projectile this becomes
\begin{align}\label{eq:JIMWLKsinglet}
    \log\frac{k^+}{\Lambda}\,\frac{d^{3}N_{\rm JIMWLK}}{d^{3}k} & = \frac{1}{(2\pi)^{3}}\,\frac{g^{4}}{16\pi^{5}}\,\frac{1}{k^{+}}\,\frac{1}{N^{2}_{c}-1}\,\log{\left(\frac{k^{+}}{\Lambda}\right)}\,\int_{x,y,z,w,w'}e^{-ik\cdot(w-w')}\notag\\&\times\langle\rho^{f}(x)\rho^{f}(y) \rangle\,\Bigg[\notag \\ & \mathcal K(x-w)\cdot\mathcal K(y-w')\,K(w,w',z)\Big[4 N_{c}{\rm Tr}\{U^\dagger(w)U(w')\}\notag\\
&-2{\rm Tr}\{T^eU^\dagger(w)U(w')T^eU^\dagger(w')U(z)\}\notag\\
&\quad-2{\rm Tr}\{T^eU^\dagger(w)U(w')T^eU^\dagger(z)U(w)\}\Big]\notag\\
&+\mathcal K(x-w)\cdot\mathcal K(y-w')\,K(x,w',z)\Big[-4 N_{c}{\rm Tr}\{U^\dagger(x)U(w')\}\notag\\
&+2{\rm Tr}\{T^eU^\dagger(w')U(x)T^eU^\dagger(x)U(z)\}\notag\\
&\quad+2{\rm Tr}\{T^eU^\dagger(w')U(x)T^eU^\dagger(z)U(w')\}\Big]\notag\\
&+\mathcal K(x-w)\cdot\mathcal K(y-w')\,K(x,y,z)\Big[4 N_{c}{\rm Tr}\{U^\dagger(x)U(y)\}\notag\\
&-2{\rm Tr}\{T^eU^\dagger(x)U(y)T^eU^\dagger(y)U(z)\}\notag\\
&\quad-2{\rm Tr}\{T^eU^\dagger(x)U(y)T^eU^\dagger(z)U(x)\}\Big]\notag\\
&+\mathcal K(x-w)\cdot\mathcal K(y-w')\,K(y,w,z)\Big[-4 N_{c}{\rm Tr}\{U^\dagger(w)U(y)\}\notag\\
&+2{\rm Tr}\{T^eU^\dagger(w)U(y)T^eU^\dagger(y)U(z)\}\notag\\
&\quad+2{\rm Tr}\{T^eU^\dagger(w)U(y)T^eU^\dagger(z)U(w)\}\Big]\notag\\
&+\mathcal K(x-w)\cdot\mathcal K(y-w')\,K(w',w',z)\Big[\notag\\
&\quad-2 N_{c}{\rm Tr}\{U^\dagger(w)U(w')\}+2 N_{c}{\rm Tr}\{U^\dagger(x)U(w')\}\notag\\
&-{\rm Tr}\{T^eU^\dagger(w')U(x)T^eU^\dagger(w')U(z)\}\notag\\
&\quad-{\rm Tr}\{T^eU^\dagger(w')U(x)T^eU^\dagger(z)U(w')\}\notag\\
&+{\rm Tr}\{T^eU^\dagger(w)U(w')T^eU^\dagger(w')U(z)\}\notag\\
&\quad+{\rm Tr}\{T^eU^\dagger(w)U(w')T^eU^\dagger(z)U(w')\}\notag\\
&+{\rm Tr}\{T^eU^\dagger(w')U(x)\}\,{\rm Tr}\{T^eU^\dagger(w')U(z)\}\notag\\
&\quad+{\rm Tr}\{T^eU^\dagger(w)U(w')\}\,{\rm Tr}\{T^eU^\dagger(w')U(z)\}\Big]\notag\\
&+\mathcal K(x-w)\cdot\mathcal K(y-w')\,K(w,w,z)\Big[\notag\\
&\quad-2 N_{c}{\rm Tr}\{U^\dagger(w)U(w')\}+2 N_{c}{\rm Tr}\{U^\dagger(w)U(y)\}\notag\\
&+{\rm Tr}\{T^eU^\dagger(w)U(w')T^eU^\dagger(w)U(z)\}\notag\\
&\quad+{\rm Tr}\{T^eU^\dagger(w)U(w')T^eU^\dagger(z)U(w)\}\notag\\
&-{\rm Tr}\{T^eU^\dagger(w)U(y)T^eU^\dagger(w)U(z)\}\notag\\
&\quad-{\rm Tr}\{T^eU^\dagger(w)U(y)T^eU^\dagger(z)U(w)\}\notag\\
&-{\rm Tr}\{T^eU^\dagger(w)U(w')\}\,{\rm Tr}\{T^eU^\dagger(w)U(z)\}\notag\\
&\quad+{\rm Tr}\{T^eU^\dagger(w)U(y)\}\,{\rm Tr}\{T^eU^\dagger(w)U(z)\}\Big]\notag\\
&+\mathcal K(x-w)\cdot\mathcal K(y-w')\,K(x,x,z)\Big[\notag\\
&\quad2 N_{c}{\rm Tr}\{U^\dagger(x)U(w')\}-2 N_{c}{\rm Tr}\{U^\dagger(x)U(y)\}\notag\\
&-{\rm Tr}\{T^eU^\dagger(w')U(x)T^eU^\dagger(x)U(z)\}\notag\\
&\quad-{\rm Tr}\{T^eU^\dagger(w')U(x)T^eU^\dagger(z)U(x)\}\notag\\
&+{\rm Tr}\{T^eU^\dagger(x)U(y)T^eU^\dagger(x)U(z)\}\notag\\
&\quad+{\rm Tr}\{T^eU^\dagger(x)U(y)T^eU^\dagger(z)U(x)\}\notag\\
&-{\rm Tr}\{T^eU^\dagger(w')U(x)\}\,{\rm Tr}\{T^eU^\dagger(x)U(z)\}\notag\\
&\quad-{\rm Tr}\{T^eU^\dagger(x)U(y)\}\,{\rm Tr}\{T^eU^\dagger(x)U(z)\}\Big]\notag\\
&+\mathcal K(x-w)\cdot\mathcal K(y-w')\,K(y,y,z)\Big[\notag\\
&\quad2 N_{c}{\rm Tr}\{U^\dagger(w)U(y)\}-2 N_{c}{\rm Tr}\{U^\dagger(x)U(y)\}\notag\\
&-{\rm Tr}\{T^eU^\dagger(w)U(y)T^eU^\dagger(y)U(z)\}\notag\\
&\quad-{\rm Tr}\{T^eU^\dagger(w)U(y)T^eU^\dagger(z)U(y)\}\notag\\
&+{\rm Tr}\{T^eU^\dagger(x)U(y)T^eU^\dagger(y)U(z)\}\notag\\
&\quad+{\rm Tr}\{T^eU^\dagger(x)U(y)T^eU^\dagger(z)U(y)\}\notag\\
&-{\rm Tr}\{T^eU^\dagger(w)U(y)\}\,{\rm Tr}\{T^eU^\dagger(y)U(z)\}\notag\\
&\quad+{\rm Tr}\{T^eU^\dagger(x)U(y)\}\,{\rm Tr}\{T^eU^\dagger(y)U(z)\}\Big]\Bigg]
\end{align} 
where the products of traces with $T^e$ are the adjoint rewriting of the $ffUUU$ contractions of Eq.~\eqref{eq:JIMWLKrho}.

Eq.~\eqref{eq:JIMWLKrho} is the JIMWLK evolution of the target operator of the LO spectrum over the rapidity interval $\log(k^+/\Lambda)$. With the JIMWLK Hamiltonian in the normalisation of~\cite{Lublinsky:2016meo},
\begin{align}\label{eq:HJIMWLK}
H_{\rm JIMWLK}=\frac{\alpha_s}{2\pi^2}\int_{x,y,z}K(\v x,\v y,\v z)\Big[J_L^e(\v x)J_L^e(\v y)+J_R^e(\v x)J_R^e(\v y)-2\,c^{ef}(\v z)\,J_L^e(\v x)J_R^f(\v y)\Big]\,,
\end{align}
where $J_L^e(\v x)$ and $J_R^e(\v x)$ act on a Wilson line at $\v x$ as $J_L^eU=T^eU$ and $J_R^eU=UT^e$, $c^{ef}(\v z)$ is defined by $U^\dagger(\v z)T^eU(\v z)=c^{ef}(\v z)T^f$, and $\frac{d}{dY}=-H_{\rm JIMWLK}$, we find
\begin{align}\label{eq:JIMWLKasLO}
\log\frac{k^+}{\Lambda}\,\frac{d^3N_{\rm JIMWLK}}{d^3k}=-\log\frac{k^+}{\Lambda}\;H_{\rm JIMWLK}\,\frac{d^3N_{\rm LO}}{d^3k}\,,
\end{align}
where $H_{\rm JIMWLK}$ acts on the Wilson lines of Eq.~\eqref{eq:LO} only. The prefactors match since $\frac{g^4}{16\pi^5}\big/\frac{g^2}{2\pi^2}=\frac{\alpha_s}{2\pi^2}$. Because $\mathcal O_{\rm LO}$ is bilinear in $U$, $H_{\rm JIMWLK}$ acting on it can be evaluated exactly, for an arbitrary (not necessarily colour-singlet or symmetric) ordered correlator $\langle\rho^{b'}(\v x')\rho^b(\v x)\rangle$; Eq.~\eqref{eq:JIMWLKrho} agrees with the result pointwise in all transverse coordinates, and Eq.~\eqref{eq:JIMWLKsinglet} with its colour-singlet projection (Section~\ref{sec:checks}). For a single adjoint dipole the action of $H_{\rm JIMWLK}$ reduces to the BK form, $-H_{\rm JIMWLK}{\rm Tr}[U^\dagger(\v a)U(\v b)]=\frac{\alpha_s}{\pi^2}\int_zK_D(\v a,\v b,\v z)\big[c^{ef}(\v z){\rm Tr}(U^\dagger(\v a)T^eU(\v b)T^f)-N_c{\rm Tr}(U^\dagger(\v a)U(\v b))\big]$.

\subsection{Why the logarithms are the evolution of the LO spectrum}\label{sec:proofsketch}

The identifications~\eqref{eq:BFKLasLO} and~\eqref{eq:JIMWLKasLO} are the expansion to first order of the leading-logarithmic factorisation of the single inclusive spectrum~\cite{Kovner:2006wr,Kovner:2006ge,Gelis:2008rw}: the spectrum of a gluon at rapidity $y$ is the LO expression evaluated with the projectile evolved from its own rapidity $Y_P$ down to $y$ and with the target evolved from $y$ to its own rapidity,
\begin{align}\label{eq:LLfact}
\frac{d^3N}{d^3k}\bigg|_{\rm LL}=\Big\langle\,\frac{d^3N_{\rm LO}}{d^3k}\,\Big\rangle_{P,\,Y_P-y;\;T,\,y-Y_T}\,.
\end{align}
The structure of the calculation shows why this holds, and we sketch the argument. Consider first the region $p^+\gg k^+$. Here the second gluon is much faster than the observed one, and its emission and absorption amplitudes factorise from those of the observed gluon: in all coefficients the dependence on $k^+$ disappears at leading power in $k^+/p^+$, and the observed gluon is emitted by the total colour charge of the valence charges and of the fast gluon, $\rho^a(\v x)\to\rho^a(\v x)+$ (charge density of the gluon at $\v z$). The sum over the rows is then the LO spectrum evaluated in the projectile state that contains one more soft gluon, minus the probability that it is not there (the virtual terms with $\mf N_2$ and the self-energy parts of $\mf A^{(3)}$), i.e.\ the LO spectrum with the projectile correlator evolved by one step. For a dilute projectile this step is linear in $\langle\rho\rho\rangle$, Eq.~\eqref{eq:BFKLrhorho}; terms with three and four charges, which would describe the non-linear (KLWMIJ) evolution of a dense projectile, must and do cancel at this order in the symmetrised sense (Sections~\ref{sec:fourrho},~\ref{sec:threerho}). In the region $p^+\ll k^+$ the second gluon is much slower than the observed one. It is emitted coherently by the valence charges and by the observed gluon, and at leading power its emission does not change the momentum of the observed gluon. The rows then combine into the change of the scattering amplitudes of the valence charges and of the observed gluon due to one soft gluon that scatters together with them, i.e.\ into one step of the JIMWLK evolution of the Wilson lines of the LO amplitude and of its conjugate, including the virtual terms. Both statements are verified by the explicit results~\eqref{eq:BFKLrho}--\eqref{eq:JIMWLKsinglet}.

Two features of the calculation enter this argument in an essential way. First, the logarithms are split at the rapidity of the observed gluon, Eq.~\eqref{eq:logsplit}; the four-$\rho$ terms and the normalisation, which carry $\log(\vee/\Lambda)$, are distributed between the two evolutions accordingly. Second, the operator ordering of the charges matters: it is the commutator terms of the reordering that produce part of the two-$\rho$ structures of both evolutions.

\subsection{Summary of the logarithmic terms}\label{sec:logsummary}

The rapidity-enhanced part of the $O(g^4)$ cross section is
\begin{align}\label{eq:logsummary}
\frac{d^3N}{d^3k}\bigg|_{\log}=\log\frac{\vee}{k^+}\,\frac{d^3N_{\rm BFKL}}{d^3k}+\log\frac{k^+}{\Lambda}\,\frac{d^3N_{\rm JIMWLK}}{d^3k}\,,
\end{align}
with the BFKL term~\eqref{eq:BFKLrho} and the JIMWLK term~\eqref{eq:JIMWLKrho}, and no term with three or four charges multiplies a rapidity logarithm. Individual rows contain also double logarithms, $\log^2(k^+/\Lambda)$, from the region in which the second gluon is both soft and far from the observed one in the transverse plane. They come from the recoil of the observed gluon in the rows with $\mf B_3$ and $\mf C_1$ and cancel in the sum over rows (Section~\ref{sec:fin_soft}).
% Source of the ln^2 cancellation: limits_checked/Cancellation_of_flagged_terms.tex (numerical, SU(2), SU(3)), presented in sec_finite.
